\section*{Capítulo \ref{ch:b3:petrological-thermodynamics} --- Termodinámica de las rocas}

\begin{solution}{exo:b3:petrological-thermodynamics:1}
$\Delta_rH = -3029.3 - 910.7 + 3935.0 = \qty{-5.0}{kJ}$; $\Delta_rS = 133.5 +
41.5 - 207.4 = \qty{-32.4}{J.K^{-1}}$; $\Delta_rV = 60.40 + 22.69 - 100.07 =
\qty{-16.98}{cm^3} = \qty{-1.698}{J.bar^{-1}}$. Pendiente $32.4/1.698 =
\qty{19.1}{bar/K}$. A \qty{873}{K}: $P = 1 + (873.15\times(-32.4) +
5000)/(-1.698) = \qty{13.7}{kbar}$.
\end{solution}

\begin{solution}{exo:b3:petrological-thermodynamics:2}
$c = 2$ y $p = 3 > c$: la regla no se cumple. Las tres fases coexisten solo en
la línea univariante de \ce{Mg2SiO4 + SiO2 <=> 2 MgSiO3}; una roca escogida
al azar no está allí, así que la asociación registra una reacción incompleta, por
ejemplo, forsterita recubierta por un borde de enstatita que la separó del
cuarzo.
\end{solution}

\begin{solution}{exo:b3:petrological-thermodynamics:3}
$a_{\text{grosularia}} = 0.20^3 = 0.008$; $a_{\text{pirope}} = 0.15^3 =
\num{3.4e-3}$; $a_{\text{anortita}} = 0.40$.
\end{solution}

\begin{solution}{exo:b3:petrological-thermodynamics:4}
$T_c = 15\,000/(2\times8.314) = \qty{902}{K}$; $W = 2\times8.314\times1000
= \qty{16.6}{kJ.mol^{-1}}$.
\end{solution}

\begin{solution}{exo:b3:petrological-thermodynamics:5}
$\Delta_rG = -5432.2 + 2\times(-2088.8) - 730.2 + 3\times3478.1 =
\qty{94.3}{kJ}$; $P = 1 + 94\,300/6.61 = \qty{14.3}{kbar}$. La
recta experimental da $22.80\times726.85 - 1093 = \qty{15.5}{kbar}$, \qty{1.2}{kbar}
superior.
\end{solution}

\begin{solution}{exo:b3:petrological-thermodynamics:6}
$RT = \qty{7483}{J.mol^{-1}}$. Para $x = 0.5$: $7483\ln0.5 + 20\,000/4 = -5187 +
5000 = \qty{-187}{J.mol^{-1}}$. Para $x = 0.111$: $7483\times(0.111\ln0.111 +
0.889\ln0.889) + 20\,000\times0.111\times0.889 = -2609 + 1974 =
\qty{-635}{J.mol^{-1}}$, y lo mismo para 0.889 por simetría. Cantidades iguales de
los dos flancos tienen la composición global 0.5 y la energía de Gibbs
\qty{-635}{J.mol^{-1}}: la separación la reduce en \qty{448}{J.mol^{-1}}.
\end{solution}

\begin{solution}{exo:b3:petrological-thermodynamics:7}
$c + 2 = 5$ fases y 5 líneas. (D) $A + C = 2B$: $0 + 2/3 = 2\times1/3$. (A)
$B + D = 2C$: $1/3 + 1 = 2\times2/3$. (C) $2A + D = 3B$: $0 + 1 = 3\times1/3$.
(B) $A + 2D = 3C$: $0 + 2 = 3\times2/3$.
\end{solution}

\begin{solution}{exo:b3:petrological-thermodynamics:8}
Por la regla de la palanca, la fase en 0.889 supone $(0.3 - 0.111)/(0.889 - 0.111) =
0.243$, \qty{24}{\percent}, y la fase en 0.111, \qty{76}{\percent}.
\end{solution}

\begin{solution}{exo:b3:petrological-thermodynamics:9}
$\partial P/\partial X_{\text{An}} = -3RT/(6.21X_{\text{An}}) =
-3\times8.314\times900/(6.21\times0.30) = \qty{-12000}{bar}$ por unidad de
$X_{\text{An}}$: un error de 0.01 desplaza la presión en \qty{120}{bar}.
\end{solution}

\begin{solution}{exo:b3:petrological-thermodynamics:10}
$G'_{\text{mix}} = RT\ln[x/(1 - x)] + W(1 - 2x)$ y $G''_{\text{mix}} =
RT/[x(1 - x)] - 2W$, cero donde $x(1 - x) = RT/(2W)$. Para $x = 0.3$: $T =
2\times20\,000\times0.21/8.314 = \qty{1010}{K}$; el \omterm{def:b3:petrological-thermodynamics:solvus}{solvus} se alcanza a
$20\,000\times0.4/(8.314\ln(7/3)) = \qty{1136}{K}$. Entre ambos valores, el
cristal es metaestable: la separación bajaría su energía, pero una pequeña
fluctuación la eleva, así que la segunda fase debe nuclear.
\end{solution}

\begin{solution}{exo:b3:petrological-thermodynamics:11}
Como $G_{\text{mix}}(x) = G_{\text{mix}}(1 - x)$, la tangente común toca
la curva en dos puntos simétricos $x$ y $1 - x$ y es horizontal:
$G'_{\text{mix}}(x) = 0$, que es la ecuación. Resolviendo, $T(x) = W(1 -
2x)/[R\ln((1 - x)/x)]$. Con $x = 1/2 - \varepsilon$, $1 - 2x = 2\varepsilon$
y $\ln((1/2 + \varepsilon)/(1/2 - \varepsilon)) \approx 4\varepsilon$, así que
$T \to W/(2R)$ cuando $\varepsilon \to 0$; y para $x < 1/2$, $T(x)$ crece a medida que $x$
se aproxima a 1/2, así que este límite es el máximo.
\end{solution}

\begin{solution}{exo:b3:petrological-thermodynamics:12}
$(2.0 - 0.08\times10)/0.92 = \qty{1.30}{wt\percent}$. El manganeso favorece
el granate: con menos de él en la \omterm{def:b3:petrological-thermodynamics:pseudosection}{composición global efectiva}, los bordes crecen en un
campo distinto y a lo largo de isopletas distintas de las que predeciría el
análisis de roca total.
\end{solution}

\begin{solution}{pb:b3:petrological-thermodynamics:1}
\textbf{1.} $\Delta_rH^\circ = -6640.0 + 2\times(-2593.8) - 910.7 + 3\times4234.0
= \qty{-36.3}{kJ}$; $\Delta_rS^\circ = 260.1 + 165.6 + 41.5 - 597.9 =
\qty{-130.7}{J.K^{-1}}$; $\Delta_rV = 125.28 + 88.30 + 22.69 - 302.37 =
\qty{-66.10}{cm^3} = \qty{-6.61}{J.bar^{-1}}$.

\textbf{2.} $-36\,300 + 298.15\times130.7 = \qty{+2.67}{kJ}$: positivo, así que
la anortita es estable en la superficie.

\textbf{3.} $P = 1 + (900\times(-130.7) + 36\,300)/(-6.61) = \qty{12.3}{kbar}$;
pendiente $130.7/6.61 = \qty{19.8}{bar/K}$.

\textbf{4.} $1 + 81\,500/6.61 = \qty{12.3}{kbar}$ a \qty{900}{K}; $1 +
119\,800/6.61 = \qty{18.1}{kbar}$ a \qty{1200}{K}.

\textbf{5.} A 626.85 y \qty{926.85}{\celsius} la recta experimental da
13.2 y \qty{20.0}{kbar}: 0.9 y \qty{1.9}{kbar} por encima de las tablas.

\textbf{6.} $1914\times6.61 = \qty{12.7}{kJ}$. La incertidumbre combinada es
$\sqrt{3.2^2 + (3\times2.0)^2 + (2\times2.0)^2 + 1.0^2} = \qty{8.0}{kJ}$: la
diferencia es 1.6 veces esa incertidumbre, algo posible en calorimetría, y la razón de que la recta
se tome de los experimentos.

\textbf{7.} $\Delta_rC_p$ es pequeño entre los sólidos, y sus dos integrales
se compensan en parte: las tablas desplazan la línea solo 26 y \qty{112}{bar}.

\textbf{8.} \ce{SiO2}, \ce{Al2O3}, \ce{FeO}, \ce{MgO}, \ce{MnO}, \ce{CaO},
\ce{Na2O}, \ce{K2O} y \ce{H2O}: $c = 9$. Fases: granate, biotita,
moscovita, plagioclasa, cianita, cuarzo y fluido, $p = 7$.

\textbf{9.} $F = 9 - 7 + 2 = 4$: presión, temperatura y dos
variables composicionales. $p = 7 \le c = 9$: la regla se cumple, y la asociación puede ser estable
en una región del plano presión--temperatura.

\textbf{10.} Dos polimorfos puros de \ce{Al2SiO5} en equilibrio se sitúan sobre la
línea univariante cianita = silimanita: una relación entre $P$ y $T$, y un
termómetro proporciona entonces la presión.

\textbf{11.} El fluido impone el potencial químico del agua: 8 componentes libres
para 6 fases sólidas, $F = 8 - 6 + 2 = 4$, y $p \le c$ se sigue
cumpliendo.

\textbf{12.} $a_{\text{grosularia}} = 0.10^3 = 0.001$, $a_{\text{anortita}} =
0.30$; $\ln K = \ln(0.001/0.30^3) = -3.30$.

\textbf{13.} $K_D = (0.12/0.74)/(0.45/0.55) = 0.1622/0.8182 = 0.198$.

\textbf{14.} $T = (52\,635 + 2488)/(20.20 + 24.94\times1.618) = 55\,123/60.57
= \qty{910}{K}$, o \qty{637}{\celsius}.

\textbf{15.} $P_0 = 22.80\times636.9 - 1093 = \qty{13429}{bar}$; $(3\times8.314
\times910.1/6.21)\ln(1/3) = \qty{-4016}{bar}$; $P = \qty{9.41}{kbar}$.

\textbf{16.} A \qty{9.41}{kbar}, \qty{644.2}{\celsius}, y luego
\qty{9.55}{kbar}; a esa presión, \qty{644.9}{\celsius} y
\qty{9.56}{kbar}, sin cambios después: \qty{645}{\celsius} y
\qty{9.56}{kbar}.

\textbf{17.} \qty{641}{\celsius} y \qty{8.79}{kbar}. Las tablas sitúan la
línea de miembro extremo unos \qty{1}{kbar} demasiado baja (Parte I); la línea experimental
se conserva.

\textbf{18.} Con $X_{\text{an}} = 0.35$: \qty{642}{\celsius} y
\qty{8.93}{kbar}, \qty{0.63}{kbar} menos.

\textbf{19.} A \qty{918}{K}: $P = 1 + (918\times12.6 - 7700)/0.571 =
\qty{6.8}{kbar}$, por debajo de \qty{9.56}{kbar}: cianita, como se observa.

\textbf{20.} $z = \num{9.56e8}/(2800\times9.81) = \qty{34.8}{km}$.

\textbf{21.} $\Delta P = (3\times8.314\times918/6.21)\ln(0.07/0.10) =
\qty{-1.3}{kbar}$: el núcleo creció unos \qty{4.8}{km} más someros. El granate
creció durante el enterramiento, y su borde registra el punto más profundo.

\textbf{22.} El esquisto de Varden alcanzó una presión máxima de
\textbf{\qty[text-series-to-math=true]{9.6}{kbar}} a
\qty{645}{\celsius}, con un enterramiento de unos \qty{35}{km}.
\end{solution}
